\documentclass[11pt]{article}
\usepackage{fiche}
\matiere{Support Vecteur Machine}
\chapitre{M2}
\source{SVM.pdf}

\begin{document}
%% page: SVM.pdf (p. 3)
\titre{\couleur{rose}{\entoure{2} \souligne{Support vecteur machine}}}

\begin{itemize}
\item Sous la forme dual, le programme SVM s'écrit
\begin{encadre}
\[ \lambda^* = \underset{\lambda_1, \dots, \lambda_N}{\arg\max} \sum_{i=1}^{N} \lambda_i - \frac{1}{2} \sum_{i=1}^{N} \sum_{j=1}^{N} \lambda_i \lambda_j y_i y_j \langle x_i, x_j \rangle \]
\end{encadre}
\couleur{rouge}{\entoure{sc : $\lambda_i \geq 0$ \\ sc : $\sum_{i=1}^{N} \lambda_i y_i = 0$}}
\end{itemize}

\marge{\couleur{vert}{\souligne{NB} : \\ La matrice de \souligne{gram} \\ $\underset{(N,N)}{\Omega} = \begin{pmatrix} \langle x_1, x_1 \rangle & \cdots & \langle x_1, x_N \rangle \\ \vdots & & \vdots \\ \langle x_N, x_1 \rangle & \cdots & \langle x_N, x_N \rangle \end{pmatrix}$}}

\couleur{rouge}{La solut° optimale de \souligne{SVM}}
\[ w^* = \sum_{i=1}^{N} \lambda_i^* y_i x_i = \sum_{i \in S} \lambda_i^* y_i x_i \quad \left\{ \begin{array}{l} \text{vu que tous les } \lambda_i^* = 0 \\ \text{pour les } i \notin \text{vecteur de support } (S) \end{array} \right. \]
\[ b^* = y_i - \langle w^*, x_i \rangle, \ \forall i \in S \quad \left\{ \begin{array}{l} \text{C'est vrai, car :} \\ y_i [\langle w^*, x_i \rangle + b^*] = 1 \Rightarrow b^* = y_i - \langle w^*, x_i \rangle \end{array} \right. \]

\couleur{rouge}{$\hookrightarrow$ On s'intéresse qu'aux points proche de l'hyperplan, le reste on s'en fout \\ $\to$ Le programme s'optimise sur les vecteurs de support et c'est tout !}

\dessin[3,26,16,36]{croix rouges à gauche et ronds bleus à droite, droite séparatrice verticale avec zone hachurée passant par les points les plus proches}

\marge{\couleur{vert}{La droite séparatrice entre les classes distinctes dépend que de \illisible{}, le reste, tout on \souligne{s'en fout}}}

\[ h(x) = \langle w^*, x \rangle + b^* = \sum_{i \in S} \lambda_i y_i \langle x_i, x \rangle + b^* \]

\couleur{rose}{Primal}
\[ (w^*, b^*) = \underset{w \in \mathbb{R}^d, b \in \mathbb{R}}{\arg\min} \frac{1}{2} \|w\|^2 \quad \text{sc : } y_i h(x_i) \geq 1, \ \forall i = 1, \dots, N \]
\[ h(x) = \langle w^*, x \rangle + b^* \]

\couleur{rose}{Dual}
\[ \lambda^* = \arg\max \sum_{i=1}^{N} \lambda_i - \frac{1}{2} \sum_{i=1}^{N} \sum_{j=1}^{N} \lambda_i \lambda_j y_i y_j \langle x_i, x_j \rangle \quad \text{sc : } \lambda_i \geq 0, \ \text{sc : } \sum_{i=1}^{N} \lambda_i y_i = 0 \]
\[ w^* = \sum_{i \in S} \lambda_i^* y_i x_i \qquad b^* = y_i - \langle w^*, x_i \rangle \]
\[ h(x) = \sum_{i \in S} \lambda_i^* y_i \langle x_i, x \rangle + b^* \]

\begin{itemize}
\item Avantage et limites du \souligne{primal/dual}
\begin{enumerate}
\item Solut° primal = Solut° dual
\item Vecteurs support classe toutes les obs en \souligne{dual}
\item $\langle x_i, x_j \rangle$ introduit le truc Kernel
\item \couleur{rouge}{Si $d \gg n$, dual est \souligne{meilleur} \\ primal = opti en $d+1$ condit° ($w \in \mathbb{R}^d$ et $b \in \mathbb{R}$) \\ dual = optimisat° sur $n$ condit° ($\lambda_1, \dots, \lambda_N$)}
\end{enumerate}
\end{itemize}

\begin{itemize}
\item Et si, l'échantillon n'est pas parfaitement linéaire séparable ?

\couleur{rouge}{\entoure{2} cas : (i) L'échant est \souligne{presque} séparable, mais quelques obs ne peuvent pas être correctement classé \\ (ii) L'échantillon admet une séparation optimale \souligne{non-linéaire}}
\end{itemize}
\end{document}
